%%%%%%%% ICML 2025 EXAMPLE LATEX SUBMISSION FILE %%%%%%%%%%%%%%%%%

\documentclass{article}

% Recommended, but optional, packages for figures and better typesetting:
\usepackage{microtype}
\usepackage{graphicx}
\usepackage{subfigure}
\usepackage{booktabs} % for professional tables
% \usepackage{array}

% hyperref makes hyperlinks in the resulting PDF.
\usepackage{hyperref}

\usepackage{capt-of}

% Attempt to make hyperref and algorithmic work together better:
\newcommand{\theHalgorithm}{\arabic{algorithm}}

% Use the following line for the initial blind version submitted for review:
\usepackage[accepted]{icml2025}
% \icmlDisableLineNumbers

% For theorems and such
\usepackage{amsmath}
\usepackage{amssymb}
\usepackage{mathtools}
\usepackage{amsthm}

% if you use cleveref..
\usepackage[capitalize,noabbrev]{cleveref}

\usepackage{subcaption} % in preamble

%%%%%%%%%%%%%%%%%%%%%%%%%%%%%%%%
% THEOREMS
%%%%%%%%%%%%%%%%%%%%%%%%%%%%%%%%
\theoremstyle{plain}
\newtheorem{theorem}{Theorem}[section]
\newtheorem{proposition}[theorem]{Proposition}
\newtheorem{lemma}[theorem]{Lemma}
\newtheorem{corollary}[theorem]{Corollary}
\theoremstyle{definition}
\newtheorem{definition}[theorem]{Definition}
\newtheorem{assumption}[theorem]{Assumption}
\theoremstyle{remark}
\newtheorem{remark}[theorem]{Remark}

\usepackage[textsize=tiny]{todonotes}

\icmltitlerunning{Explainable Clustering for Spatio-Temporal Sensor Data}

\begin{document}

\twocolumn[
\icmltitle{Literature Review:\\
``Explainable Clustering for Spatio-Temporal Sensor Data''}

\begin{icmlauthorlist}
\icmlauthor{An Pham}{equal}
\end{icmlauthorlist}

\icmlaffiliation{inst1}{Faculty of Science and Technology, Université Jean Monnet,
Saint-Étienne, France}

\icmlcorrespondingauthor{An Pham}{binh-an.pham@lirmm.fr; binh-an.pham@umontpellier.fr}

\icmlkeywords{Explainable Clustering, Spatio-Temporal Data, Interpretable Machine Learning,
Representation Learning, SHAP, LIME, ICML}

\vskip 0.3in
]

\printAffiliationsAndNotice{}

\begin{abstract}
Clustering spatio-temporal sensor data in a way that is simultaneously
accurate and interpretable requires co-designing three interdependent
components: the \emph{data representation}, the \emph{clustering strategy},
and the \emph{explanation strategy}. This review organises the literature
around these three core axes and examines how choices made at each axis
propagate to---and constrain---the others. I show that representation
determines both clustering quality and the very possibility of meaningful
explanation; that clustering strategies range from partitioning and
density-based methods to joint-optimisation frameworks that treat
explainability as a first-class objective; and that explanation strategies span post-hoc attribution (SHAP, LIME) and intrinsic, ante-hoc designs. I further analyse the key trade-offs (raw data vs.\ engineered features; global vs.\ local explanations; performance vs.\ interpretability), identify the fundamental alignment gap between sequentially optimised pipeline stages, and enumerate open research
problems specific to the spatio-temporal, autocorrelated, and potentially incomplete sensor setting of the PICOPATT project.
\end{abstract}

% ============================================================
\section{Introduction and Problem Setting}
\label{sec:intro}
% ============================================================

Clustering sensor data from environmental monitoring systems poses a
distinctive set of challenges that standard clustering theory does not
fully address. In the PICOPATT context, each spatial station records
approximately~23 meteorological variables at four time steps per day,
yielding a multivariate time series with strong spatio-temporal
dependencies. The raw observation for a single station on a single day
can be written as $\mathbf{X} \in \mathbb{R}^{T \times V}$, where $T=4$
and $V \approx 23$, before any temporal aggregation or feature engineering.

Two structural properties immediately complicate both clustering and
explanation. First, consecutive observations are \emph{not} independent:
sensor readings exhibit strong positive autocorrelation, violating the
i.i.d.\ assumption underpinning most clustering
theory~\cite{hamdi_spatiotemporal_2022}. Second, data may be \emph{incomplete}: sensor
failures or maintenance windows create irregular gaps that must be handled
prior to or within the clustering pipeline.

A further difficulty is \emph{terminological ambiguity}: ``faithful to
what?''---to the reference partition produced by a black-box algorithm,
to the underlying physical phenomenon, or to the end user's
decision-making process---is a question the literature answers
inconsistently~\cite{portugal_framework_2024}. I adopt the position that faithfulness
should ultimately refer to the \emph{real-world phenomenon}, not merely
to the output of a surrogate model.

The central research question is therefore:
\begin{quote}
  \emph{How do I align representation, clustering, and explanation so
  that the resulting descriptions are meaningful, faithful, and
  human-intelligible in a spatio-temporal sensor context?}
\end{quote}
This question moves the problem from a \emph{post-hoc task} (explaining
a result that has already been generated) to a \emph{systemic design
challenge} (ensuring the entire pipeline is coherent). The remainder of
this review addresses this challenge by decomposing it into three
interlocking axes.

% ============================================================
\section{Three-Axis Taxonomy}
\label{sec:taxonomy}
% ============================================================

Following and extending recent surveys~\cite{hu_interpretable_2025, portugal_framework_2024,hwang_xclusters_2023},
I organise the literature into three axes that correspond to where
decisions about representation, clustering, and interpretability are made.

\paragraph{Axis~A --- Data Representation.}
How are raw sensor signals encoded before any clustering is applied?
The choice between raw time series, hand-crafted statistical features,
and learned embeddings directly determines both the clustering quality
and the interpretability of downstream explanations.

\paragraph{Axis~B --- Clustering Strategy.}
Given a feature space, which algorithm is used to partition it?
Options range from classic partitioning (k-means), density-based
(DBSCAN/HDBSCAN), and graph/spectral methods, to decision-tree approaches
and modern joint-optimisation frameworks that simultaneously minimise
cluster distortion and explanation complexity.

\paragraph{Axis~C --- Explanation Strategy.}
Once a partition exists, how is it explained? Post-hoc attribution
methods (SHAP, LIME) operate on an already-fitted model, while
intrinsic, ante-hoc approaches embed interpretability into the model's
design from the outset. Hybrid strategies combine both.

The key insight is that \textbf{representation determines both clustering
quality and explainability simultaneously}: it defines the feature space
in which clusters are formed \emph{and} the vocabulary in which
explanations are expressed. Table~\ref{tab:litreview} maps principal works
to these three axes.

% ============================================================
\section{Axis A: Data Representation}
\label{sec:representation}
% ============================================================

Representation is often underestimated or treated as a preprocessing
detail, yet it is the axis on which the entire pipeline rests. For a
station-level observation $\mathbf{X} \in \mathbb{R}^{T \times V}$, the
encoding $\phi : \mathbf{X} \mapsto \mathbb{R}^d$ shapes what the
clustering algorithm can learn and what the explanation can express.

\subsection{Raw Time Series}

Using raw, unprocessed multivariate sequences preserves the maximum
amount of information and avoids the subjective choices inherent in
feature engineering~\cite{alvarez-garcia_comprehensive_2024}. However, raw representations are high-dimensional
and difficult to explain: a cluster assignment grounded in a 92-entry
matrix ($4 \times 23$) cannot easily be expressed in terms that a
meteorologist would find actionable.

\subsection{Aggregated and Engineered Features}

Statistical summaries---mean, variance, min/max per variable per
day---are directly interpretable as sensor-level quantities and
constitute the simplest form of dimensionality reduction. Shapelet-based
methods extract recurring sub-sequences that serve as interpretable
cluster descriptors~\cite{ye_time_2009}. \textsc{Time2Feat}~\cite{bonifati_time2feat_2022}
produces aggregate descriptors that are designed to be both informative
and human-readable.

The core design question is: \emph{what specific transformation is
applied, and does the resulting feature connect back to
domain-meaningful quantities?} Transparency at this stage is a
prerequisite for downstream interpretability.

\subsection{Feature Engineering Methodologies in Practice}
\label{subsec:feat_methods}

The literature presents four representative methodologies for encoding
data and explaining the resulting clusters, each making distinct choices
about how features are constructed and how explanations are derived.

\paragraph{Clust-learn ~\cite{alvarez-garcia_comprehensive_2024}.}
This framework handles mixed-type data by treating numerical and
categorical variables separately. Numerical variables are first
standardised (zero mean, unit variance) to prevent scale-driven
dominance, then compressed via \textbf{Sparse PCA} (SPCA)---a variant
of PCA that adds an $L_1$ penalty so each extracted component is a
sparse linear combination of only a few original variables. Categorical
variables are handled via \textbf{Multiple Correspondence Analysis}
(MCA). The optimal number of components is selected by the elbow method
applied to cumulative explained variance and inertia.
For explanation, cluster assignments are recast as pseudo-labels for a
classification task: an XGBoost classifier is trained on those labels and Tree SHAP values~\cite{lundberg_unified_2017, lundberg_local_2020} are extracted to deliver both local, instance-level explanations and global cluster-level feature importance.

\paragraph{Industrial Sensor Prognostics ~\cite{cohen_shapley-based_2024}.}
Applied to multivariate time-series from 18 synchronised aerospace
sensors, this approach manually extracts seven statistical time-domain
features per signal: mean, standard deviation, minimum, first quartile,
median, third quartile, and maximum. Auxiliary variables (cycle number,
flight duration) augment the feature set to 129 dimensions per cycle;
all features are then min-max normalised. The explanation layer combines
two tools: a \textbf{SHAP permutation explainer} for continuous feature
attributions, and \textbf{SkopeRules} to generate linguistic
``if-then'' decision rules characterising each failure cluster in terms
of the original sensor values. This pairing of attribution and rule
extraction illustrates the hybrid strategy described in
Section~\ref{sec:explanation}.

\paragraph{Interpretable Time-Series Representations
(Time2Feat \& Shapelets).}
\textsc{Time2Feat}~\cite{bonifati_time2feat_2022} encodes multivariate time series by
extracting both \emph{intra-signal} (per-variable statistics, trend
coefficients) and \emph{inter-signal} (cross-correlation, lag) features,
then applies Principal Feature Analysis or user-guided annotation to
select a compact, human-interpretable subset.
Shapelet-based methods~\cite{ye_time_2009} take a complementary route:
data are projected onto a \emph{similarity space} defined by short,
representative subsequences (shapelets), and the explainer---the
\textbf{Shapelet Cluster Explanation (SCE)} module of
CDPS~\cite{elamouri_constrained_2023}---outputs the specific waveform patterns most
descriptive of each cluster, enabling visual inspection by domain experts.
Both approaches anchor explanations in waveform-level quantities that
are directly meaningful to sensor operators.

\paragraph{TELL: Direct Neural Interpretable Clustering ~\cite{peng_xai_2022}.}
Instead of manual feature engineering, \textbf{TELL} (inTerpretable
nEuraL cLustering) uses a convolutional autoencoder to learn
discriminative features directly from raw data. The traditional $k$-means
objective is recast as a neural layer, so cluster centres $\Omega$ and
feature representations are jointly optimised via batch-wise stochastic
gradient descent. Interpretability is \emph{intrinsic}: the weights of
the cluster layer correspond exactly to the cluster centroids $\Omega$,
and an attention mechanism weights the contribution of each input
data point relative to those centres---no post-hoc surrogate is
required. TELL therefore occupies the intersection of Axis~A (learned
representation) and Axis~C (ante-hoc explanation).

\subsection{Learned Embeddings}

Deep or self-supervised encoders can capture complex non-linear patterns
invisible to hand-crafted features, yielding superior clustering
quality~\cite{aljalbout_clustering_2018}. The cost is an \emph{interpretability
gap}: explanations act on the transformed space~$\phi(\mathbf{X})$, not
on raw sensor readings. The chain
\[
  \mathbf{X} \;\xrightarrow{\phi}\; \phi(\mathbf{X})
  \;\xrightarrow{\text{cluster}}\; c
  \;\xrightarrow{\text{explain}}\; \text{rules}
\]
must itself be transparent: does an explanation of $\phi(\mathbf{X})$
translate back to a meaningful statement about~$\mathbf{X}$?

\subsection{Impact on Clustering Quality and Explainability}

\paragraph{Clustering quality.}
The \textbf{curse of dimensionality}~\cite{goos_surprising_2001} means that
distance functions lose discriminative power in high-dimensional spaces.
\textbf{Data vagueness}---two sensor trajectories that are geometrically
close in feature space yet represent physically distinct phenomena (e.g.,
an intensifying vs.\ dissipating temperature anomaly)---further
challenges quality. Latent space methods can mitigate both effects but
at the cost of interpretability~\cite{oord_neural_2018}.

\paragraph{Explainability.}
The \textbf{interpretability gap}~\cite{schlegel_towards_2025} widens when
features are computed by complex encoders. Research suggests selecting
or extracting \emph{interpretable features} that are (i)~informative
enough for accurate grouping and (ii)~grounded in the original sensor
units for human comprehension~\cite{hu_interpretable_2025}. This is the
alignment problem at the representation level.

% ============================================================
\section{Axis B: Clustering Strategy}
\label{sec:clustering}
% ============================================================

\subsection{Classic Partitioning and Density-Based Methods}

\textbf{K-means} assigns each observation to its nearest centroid,
yielding cluster centers expressible as feature means---directly
human-readable when features are interpretable. \textbf{DBSCAN},
\textbf{HDBSCAN}, \textbf{ST-DBSCAN}~\cite{birant_st-dbscan_2007}, and \textbf{MDST-DBSCAN}~\cite{choi_mdst-dbscan_2021} form clusters from density-connected regions, naturally
handling noise and irregular cluster shapes common in meteorological data, but produce cluster labels without a simple parametric description.
\textbf{Spectral and graph-based} methods exploit the geometry of
pairwise similarity graphs, which can encode spatial adjacency between
sensor stations, but again produce implicit cluster descriptions.

\subsection{Decision-Tree and Rule-Based Approaches}

Decision-tree-based strategies are the most prominent line of work for
explainable clustering, as they produce compact rules that are directly
interpretable as sensor-level threshold conditions.

\paragraph{Axis-aligned threshold trees.}
The \textbf{IMM} algorithm~\cite{dasgupta_explainable_nodate} fits a binary decision
tree with exactly~$k$ leaves to a reference $k$-clustering, minimising
misassignment at each leaf. \textbf{ExKMC}~\cite{frost_exkmc_2020} generalises
IMM by allowing $k'>k$ leaves, with the \emph{price of explainability}
\[
  \Delta(k,k') = \mathrm{cost}(T_{k'}) - \mathrm{cost}(C_k)
\]
quantifying the degradation in cluster quality incurred by the
explainable approximation. For $V \approx 23$ variables, axis-aligned
rules of the form ``temperature $> \theta$'' are physically meaningful
but fail to capture variable correlations.

\paragraph{Oblique and sparse linear trees.}
\textbf{TAO}~\cite{gabidolla_optimal_2022} solves the non-convex problem of fitting
oblique trees---where each split is a sparse linear combination
$\mathbf{w}^\top\phi(\mathbf{X}) > b$---by alternating between
optimising tree structure and node parameters. In a spatio-temporal
context, $\mathbf{w}$ can encode linear combinations across spatial
stations \emph{and} temporal lags, yielding semantically richer rules
at the cost of cognitive complexity.

\paragraph{Hypercubic approximations.}
\textbf{ExACT}~\cite{sabbatini_exact_nodate} partitions the feature space into
axis-aligned hypercubes, yielding interval-based rules
($\theta_{\rm lo} \le x_j \le \theta_{\rm hi}$) directly translatable
to sensor threshold alerts.

\subsection{Deep Clustering}

Deep clustering methods jointly optimise a neural encoder and a
clustering objective, achieving state-of-the-art partitions on complex
data~\cite{aljalbout_clustering_2018}. However, as discussed in
Section~\ref{sec:representation}, they compound the interpretability gap:
the resulting partition is a product of millions of gradient-based
micro-decisions, none of which is individually inspectable.

\subsection{Joint-Optimisation Frameworks}

The most principled response to the quality--interpretability
tension is to optimise both objectives simultaneously.

\textbf{XClusters}~\cite{hwang_xclusters_2023} treats explainability as a first-class
citizen:
\[
  \min_{\mathcal{C},\,T} \;\;
  \underbrace{\mathrm{distortion}(\mathcal{C})}_{\text{cluster quality}}
  + \lambda \cdot \underbrace{|T|}_{\text{tree complexity}},
\]
alternating between improving cluster assignment~$\mathcal{C}$ given the
explaining tree~$T$, and improving~$T$ given~$\mathcal{C}$.
The hyperparameter~$\lambda$ explicitly controls the tradeoff; in the
sensor context, it can be calibrated to bound the tree depth at a
domain-readable limit~$d_{\max}$.

\textbf{Algorithm-agnostic pipelines} such as Clust-learn and
MAACLI~\cite{portugal_framework_2024} provide a guided sequential pipeline---
Preprocessing~$\to$ Dimensionality Reduction~$\to$ Clustering~$\to$
Explanation---remaining agnostic to the specific algorithm at each stage.
The known limitation is that sequential optimisation does not enforce
\emph{alignment} between stages: the reduction chosen in step~2 may
discard features that would have been useful in the explanation produced in step~4. This sacrifice of joint optimisation for modularity is an active research area.
% ~\cite{TODO_joint_opt}.

\paragraph{The problem with complexity.}
Research shows that forcing a structure to be understandable increases
the within-cluster sum of squared errors compared to an unrestricted
optimal clustering~\cite{dasgupta_explainable_nodate,frost_exkmc_2020}. In high-stakes domains
(e.g., safety-critical sensor monitoring), a decision that cannot be
justified is often unusable even if it is statistically accurate; the
domain context therefore determines which side of the tradeoff to
prioritise~\cite{schlegel_towards_2025}.

% ============================================================
\section{Axis C: Explanation Strategy}
\label{sec:explanation}
% ============================================================

The fundamental tension in explanation is between \textbf{faithfulness}
(how accurately the explanation reflects the model) and \textbf{simplicity}
(how easy it is for a human to understand).

\subsection{Post-Hoc Explanations}

Post-hoc methods treat the cluster assignments from a complex algorithm
as pseudo-labels and then approximate those assignments with an
interpretable surrogate.

\paragraph{SHAP.}
\textbf{SHAP (SHapley Additive exPlanations)}~\cite{lundberg_unified_2017} assigns
to each feature~$j$ a contribution~$\phi_j$ satisfying axioms of
efficiency, symmetry, and linearity. \textbf{SHAP-C}~\cite{cohen_shapley-based_2024}
applies this to clustering by using the distance to the assigned centroid
as the value function:
\[
  g(\mathbf{z}) = \phi_0 + \textstyle\sum_{j=1}^d \phi_j z_j
  \approx -\|\phi(\mathbf{X}) - \boldsymbol{\mu}_c\|^2.
\]
A particularly interesting direction clusters instances directly in the
\textbf{SHAP explanation space}~\cite{hwang_xclusters_2023}: two sensor records are
grouped not because their raw values are similar, but because the model
assigns them similar Shapley values---i.e., they are grouped for the
\emph{same reasons}. This surfaces distinct archetypes of explanatory
behaviour.

\paragraph{LIME.}
\textbf{LIME}~\cite{ribeiro_why_2016-1} fits a sparse linear model in a
perturbed neighbourhood of each instance. It has known \emph{instability}
issues: small perturbations or changes in the sampling distribution can
yield substantially different explanations~\cite{alvarez-garcia_comprehensive_2024}. This
instability is especially problematic for autocorrelated sensor data,
where the local neighbourhood is not well-defined in the i.i.d.\ sense.
SHAP is therefore generally preferred for its stronger theoretical
guarantees, though it incurs higher computational cost.

\subsection{Intrinsic (Ante-Hoc) Strategies}

Interpretability built into the model design itself---no surrogate
needed. Table~\ref{tab:antehoc} summarises the main approaches.

\begin{table}[h]
\centering
\footnotesize
\setlength{\tabcolsep}{3pt}
\begin{tabular}{p{2.3cm}p{3.8cm}}
\toprule
\textbf{Method} & \textbf{Why it is interpretable} \\
\midrule
K-means & Cluster centers are readable feature means \\
Decision tree clustering & Splits give explicit if-then rules \\
Prototype-based & Representative points explain each cluster \\
Sparse clustering & Algorithm selects which features matter \\
Fuzzy clustering & Membership degrees show partial belonging \\
\bottomrule
\end{tabular}
\caption{Ante-hoc clustering strategies and their interpretability mechanism.}
\label{tab:antehoc}
\vspace{-1.2em}
\end{table}

\subsection{Hybrid Strategies}

Hybrid approaches combine an expressive model (for quality) with a
lightweight ante-hoc component (for local transparency). Joint
optimisation frameworks such as XClusters~\cite{hwang_xclusters_2023} and
ExKMC~\cite{frost_exkmc_2020} fall into this category: they retain a
complex clustering objective while optimising an explanatory tree
alongside it.

If you prioritise \textbf{faithfulness}, you often end up with complex
models where explanations are only approximate; if you prioritise
\textbf{simplicity}, you gain transparency but may sacrifice the model's
ability to capture complex patterns in raw sensor data.

% ============================================================
\section{Spatio-Temporal Specific Challenges}
\label{sec:spatiotemporal}
% ============================================================

Standard clustering assumes data points are i.i.d., but spatio-temporal
sensor data breaks this assumption in several ways.

\subsection{Autocorrelation and Non-Stationarity}

Sensor readings at time~$t$ are strongly correlated with readings at
time~$t{-}1$. For a univariate signal~$x_t$, the autocorrelation at
lag~$\ell$ is
\[
  \rho(\ell) = \frac{\mathrm{Cov}(x_t, x_{t-\ell})}{\mathrm{Var}(x_t)}.
\]
High~$\rho(\ell)$ for small~$\ell$ inflates effective sample size
estimates and invalidates the i.i.d.\ assumption~\cite{hamdi_spatiotemporal_2022}.
\textbf{Temporal focusing}---restricting distance computations to a
sliding time window---partially mitigates this. Non-stationarity (seasonal
distributional shift) further complicates global clustering models and
the stability of explanations over time~\cite{schlegel_towards_2025}.

\subsection{Temporal Evolution and Cluster Dynamics}

Standard clustering treats the dataset as a static snapshot, but sensor
streams are \emph{dynamic}. The literature identifies up to~14 types of
cluster transitions---birth, death, merge, split, expansion,
contraction---and proposes representing them as an \textbf{evolution
graph}~\cite{hamdi_spatiotemporal_2022}: nodes are cluster states at each time step;
directed edges encode transitions. For the PICOPATT dataset, this graph
has one node per cluster per quarter-day interval. Explaining \emph{why}
a cluster splits or merges (in terms of which variables drove the
transition) is an open problem not addressed by static explanation
methods.

\subsection{Missing Data and Irregular Sampling}

With~92 entries per station-day ($4 \times 23$), sensor malfunctions
introduce missing values. Common strategies---mean/median imputation,
low-rank matrix completion, spatial-neighbour imputation---each alter the
downstream feature distribution and hence cluster assignments, yet this
dependency is rarely made explicit in explanations.

\subsection{Semantic Ambiguity and the Right Unit of Clustering}

Two sensor trajectories may be geometrically close in feature space yet
represent physically distinct phenomena (e.g., a temperature sensor
``intensifying'' vs.\ one ``dissipating''). Standard distance-based
clustering assigns them to the same cluster, yet their physical
interpretation is opposite. Explanation methods operating purely on
feature values inherit this ambiguity~\cite{portugal_framework_2024}.

A related open question is: \emph{what is the right unit of clustering?}
Trajectories? Fixed time windows? Event episodes?~\cite{kisilevich_spatio-temporal_nodate}
The choice of temporal granularity determines what patterns are visible
and what explanations are possible.

% ============================================================
\section{The Alignment Problem: Joint Optimisation}
\label{sec:alignment}
% ============================================================

The core insight motivating joint optimisation is that the sequential
pipeline (represent $\to$ cluster $\to$ explain) introduces a fundamental
\textbf{alignment gap}: each stage is optimised independently, so the
representation may discard features useful for explanation, and the
explanation may be faithful to the clustering but not to the phenomenon.

\paragraph{The dominant paradigm and its limitation.}
The majority of explainable clustering works---including the influential
IMM~\cite{dasgupta_explainable_nodate} and ExKMC~\cite{frost_exkmc_2020}
frameworks---assume a \emph{fixed reference clustering}: the partition is
produced first by an unrestricted algorithm (e.g., $k$-means), and a
decision tree is then fitted to approximate it as a post-hoc step.
This design choice is convenient but structurally flawed: the clustering
algorithm has no incentive to produce a partition that is easy to explain,
because it is optimised without any knowledge that an explanation will
be required. The explainer is therefore forced into an inherently lossy
approximation of a partition that was never designed with interpretability
in mind.

\paragraph{From sequential fitting to co-adaptation.}
XClusters~\cite{hwang_xclusters_2023} directly addresses this by requiring the
clustering and the explaining tree to \emph{co-adapt}: neither is fixed
while the other is optimised. 
Formally, it minimises
% $\min_{\mathcal{C},\,T} D(\mathcal{C}, T) + \lambda\,N(\mathcal{C}, T)$, where $D$ measures cluster distortion and $N$ penalises tree complexity, 
$\min_{\mathcal{C},\,T} distortion(\mathcal{C}) + \lambda \cdot |T|$, where $distortion(C)$ measures cluster quality and $|T|$ penalises tree complexity, 
and the two are alternated until convergence.
Gabidolla \& Carreira-Perpi\~n\'an~\cite{gabidolla_optimal_2022} take the same
co-adaptation idea further by fitting sparse \emph{oblique} trees jointly
with the cluster assignments, so that the tree parameters and the
partition boundaries adapt to each other simultaneously.
The shared insight across both works is that explainability must be a
\emph{design constraint} on the clustering, not an afterthought applied
to its output.

\paragraph{Extension to the spatio-temporal setting.}
The objective proposed here extends this accuracy--interpretability
framework by adding two regularisers that are absent from all existing
joint-optimisation methods: a \textbf{temporal coherence} penalty
$\mathcal{L}_\mathrm{tc}$ that discourages rapid cluster switching across
consecutive time steps, and a \textbf{spatial homogeneity} penalty
$\mathcal{L}_\mathrm{sh}$ that encourages geographically adjacent stations
to share cluster assignments. For the specific case of picoclimatic sensor
data, a fully aligned joint objective would balance:
\begin{enumerate}
  \item \textbf{Temporal coherence} ($\mathcal{L}_\mathrm{tc}$): clusters
    should be stable across consecutive time steps.
  \item \textbf{Spatial homogeneity} ($\mathcal{L}_\mathrm{sh}$): nearby
    stations should tend to share clusters.
  \item \textbf{Explanation simplicity} ($|T|$): rules should involve few
    variables and thresholds.
  \item \textbf{Representational interpretability}: features used in rules
    should map back to sensor-level quantities.
\end{enumerate}
Formally, letting $\mathcal{L}_\mathrm{tc}$ and $\mathcal{L}_\mathrm{sh}$
denote the temporal and spatial penalty terms respectively:
\[
\begin{aligned}
  \min_{\mathcal{C},\,T}\;
    &\mathrm{distortion}(\mathcal{C})
     + \lambda_1\,\mathcal{L}_\mathrm{tc}(\mathcal{C}) \\
    &\quad+\;\lambda_2\,\mathcal{L}_\mathrm{sh}(\mathcal{C})
     + \lambda_3\,|T|.
\end{aligned}
\]
Current methods address at most two of these four objectives jointly;
to date, no published work optimises all four
simultaneously~\cite{portugal_framework_2024,hwang_xclusters_2023}.

\paragraph{Explaining the optimizer.}
In gradient-based clustering (deep k-means, differentiable clustering
losses), the final partition is the product of a sequence of micro-decisions
(assignment updates, centroid moves, gradient steps). SHAP values of the
final model reflect the cumulative effect of these micro-decisions but
do not reveal \emph{which} gradient steps were most consequential.
Connecting explanation methods to the optimiser via influence functions or
implicit differentiation is a promising but largely unexplored direction.


% ============================================================
\section{Trade-offs and Design Choices}
\label{sec:tradeoffs}
% ============================================================

Practitioners face several fundamental trade-offs whose resolution must
be guided by domain requirements.

% \paragraph{Raw data vs.\ engineered features.}

% \begin{table}[h]
% \centering
% \footnotesize
% \setlength{\tabcolsep}{4pt}
% \begin{tabular}{p{3cm}p{3cm}}
% \toprule
% \textbf{Raw time series} & \textbf{Engineered features} \\
% \midrule
% \checkmark\ Preserves full information & \checkmark\ Directly interpretable \\
% $\times$\ Hard to explain & $\times$\ May lose temporal structure \\
% \bottomrule
% \end{tabular}
% \caption{Raw vs.\ engineered feature trade-off.}
% \label{tab:raw_vs_feat}
% \vspace{-1.2em}
% \end{table}

\noindent\textbf{Raw data vs.\ engineered features.}\par
\vspace{0.35em}

\noindent\begin{minipage}{\columnwidth}
\centering
\footnotesize
\setlength{\tabcolsep}{4pt}
\begin{tabular}{p{3cm}p{3cm}}
\toprule
\textbf{Raw time series} & \textbf{Engineered features} \\
\midrule
\checkmark\ Preserves full information & \checkmark\ Directly interpretable \\
$\times$\ Hard to explain & $\times$\ May lose temporal structure \\
\bottomrule
\end{tabular}
\captionof{table}{Raw vs.\ engineered feature trade-off.}
\label{tab:raw_vs_feat}
\end{minipage}

\paragraph{Spatial vs.\ temporal importance.}
Are clusters driven more by station location, by temporal patterns, or
by both? This determines whether the feature encoding should emphasise
spatial adjacency (graph-based features), temporal dynamics (shapelets,
lag correlations), or multivariate cross-variable patterns.

\paragraph{Global vs.\ local explanations.}
Global (cluster-level) explanations describe the average behaviour of a
cluster; local (instance-level) explanations describe why a specific
observation was assigned to its cluster. SHAP naturally provides local
explanations that can be aggregated globally; decision trees provide
global rules but can be queried at the instance level by traversal.

\paragraph{Performance vs.\ interpretability.}
While deep learning can provide superior groupings, it often fails to
provide the human-intelligible justifications required for onboard sensor
monitoring or safety-critical tasks. In high-stakes domains, a
statistically accurate decision that cannot be justified is often
unusable. The domain context must therefore determine which objective to
prioritise, and the chosen axis must be documented and justified as part
of the system design~\cite{schlegel_towards_2025}.

% ============================================================
\section{Research Gaps and Open Problems}
\label{sec:gaps}
% ============================================================

\paragraph{Research gap.}
Current methods do not provide a \emph{consistent, structured explanation}
of clusters that integrates spatial, temporal, and sensor dimensions.
Existing explanations are: (i)~inconsistent across methods; (ii)~not
comparable across pipelines; and (iii)~not domain-aligned to the
vocabulary of the end user (e.g., meteorologist). Future methods should
be model-agnostic, scalable, and compatible with multiple clustering
approaches.

\paragraph{Stable explanations in non-stationary streams.}
How do I maintain stable, trustworthy explanations when the data
distribution shifts over time (seasonal effects, sensor
drift)?~\cite{schlegel_towards_2025} Streaming explanation methods that
update incrementally---rather than recomputing from scratch---are needed.

\paragraph{Faithfulness to the real case.}
``Faithful'' is used inconsistently: fidelity to model output, to the
training distribution, or to the underlying physical phenomenon. For
sensor data with a known generative process, faithfulness to the physical
phenomenon is the ultimate standard but is the hardest to operationalise.

\paragraph{Engineered features and their explainability.}
When features are computed by complex encoders
(e.g., Time2Feat~\cite{bonifati_time2feat_2022}), the explanation acts on the
transformed space, not on raw sensor readings. Tracing explanations back
through the transformation chain to sensor-level statements is an open
challenge.

\paragraph{Out-of-sample mapping and computational scalability.}
Most interpretable clustering methods define explanations only for
training-set instances. Mapping new, unseen sensor readings to existing
clusters and generating explanations for them without retraining remains
an underexplored practical problem, especially under real-time
monitoring requirements.

\paragraph{Summary of key references.}
Table~\ref{tab:litreview} maps the principal works surveyed to the three
axes and to relevant challenges in the PICOPATT setting.

\begin{table}[h]
\centering
\footnotesize
\setlength{\tabcolsep}{3pt}
\begin{tabular}{p{2.6cm} p{1.6cm} p{3.4cm}}
\toprule
\textbf{Reference} & \textbf{Axis} & \textbf{Contribution} \\
\midrule
\citet{dasgupta_explainable_nodate} & B (tree) & IMM; price of explainability \\
\citet{frost_exkmc_2020} & B (tree) & ExKMC; $k' > k$ leaves \\
\citet{gabidolla_optimal_2022} & B (tree) & TAO; sparse linear splits \\
\citet{sabbatini_exact_nodate} & B (tree) & ExACT; interval rules \\
\citet{lundberg_unified_2017} & C (post-hoc) & SHAP; Shapley values \\
\citet{cohen_shapley-based_2024} & C (post-hoc) & SHAP-C; centroid distance \\
\citet{ribeiro_why_2016-1} & C (post-hoc) & LIME; local surrogates \\
\citet{alvarez-garcia_comprehensive_2024} & A+C & Clust-learn; SPCA + Tree SHAP \\
\citet{hu_interpretable_2025} & A+B+C (survey) & Interpretable Clustering: A Survey \\
\citet{hwang_xclusters_2023} & B+C (joint) & XClusters; co-adapt. optim. \\
\citet{hwang_xclusters_2023} & C (SHAP) & Clustering in SHAP space \\
\citet{hamdi_spatiotemporal_2022} & ST & Evolution graphs; autocorrelation \\
\citet{portugal_framework_2024} & A+B+C & Pipeline taxonomy; MAACLI \\
\citet{bonifati_time2feat_2022} & A & Time2Feat; TS features \\
\citet{ye_time_2009} & A & Shapelet-based features \\
\citet{cohen_shapley-based_2024} & A+C & Statistical features; SkopeRules \\
\citet{elamouri_constrained_2023} & A+C & Shapelet cluster explanation \\
\citet{peng_xai_2022} & A+C (ante-hoc) & Neural clustering; attention weights \\
\bottomrule
\end{tabular}
\caption{Key references mapped to the three axes (A=Representation, B=Clustering, C=Explanation, ST=Spatio-Temporal).}
\label{tab:litreview}
\vspace{-1.5em}
\end{table}

% ============================================================
\section{Conclusion}
\label{sec:conclusion}
% ============================================================

This review has organised the literature on explainable clustering around
three core axes---data representation, clustering strategy, and
explanation strategy---and examined each axis in the context of
spatio-temporal picoclimatic sensor data.

The key takeaway is that representation is the load-bearing axis: it
determines both the clustering quality and the very vocabulary in which
explanations can be expressed. Sequential pipelines that optimise each
axis independently accumulate an alignment gap; joint-optimisation
frameworks (XClusters~\cite{hwang_xclusters_2023}, ExKMC~\cite{frost_exkmc_2020}) begin to
close this gap but have not yet been adapted to the full spatio-temporal
setting.

Open problems include: (i)~maintaining stable explanations in
non-stationary sensor streams; (ii)~defining and operationalising
faithfulness to the physical phenomenon rather than to the model;
(iii)~disambiguating semantically ambiguous trajectories through
explanation; (iv)~tracing engineered-feature explanations back to
raw sensor units; and (v)~connecting explanation methods to the
optimisation dynamics that produce the partition.
Addressing these challenges---under the unified framing of aligning
representation, clustering, and explanation---is the central goal of the
PICOPATT project.

%-------------------------------------------------------------------
% CITATION NOTES FOR MISSING REFERENCES
%-------------------------------------------------------------------
% The following \cite keys are used in the text but require .bib entries:
%
%   \cite{TODO_deep_clustering}  -- a deep clustering survey/paper
%                                   (e.g., Min et al. 2018, arXiv:1801.07648)
%   \cite{TODO_curse_dim}        -- curse of dimensionality in clustering
%                                   (e.g., Aggarwal et al. 2001, VLDB)
%   \cite{TODO_latent_space}     -- latent space clustering quality
%                                   (e.g., van den Oord et al. VQ-VAE)
%   \cite{TODO_nonstationary}    -- non-stationary / streaming XAI
%                                   (arXiv: 2507.20840)
%   \cite{TODO_alignment}        -- interpretable feature alignment
%                                   (arXiv: 2409.00743)
%   \cite{TODO_joint_opt}        -- joint optimisation / hot topic papers
%   \cite{TODO_kimana}           -- unit of clustering in trajectories
%                                   (KiMaNa10 reference)
%   \cite{TODO_tell}             -- TELL: inTerpretable nEuraL cLustering
%-------------------------------------------------------------------

\nocite{alvarez-garcia_comprehensive_2024,
        zeghina_deep_2024,
        % ribeiro_why_2016,
        peng_xai_2022,
        wang_deep_2019,
        frost_exkmc_2020,
        cohen_shapley-based_2024,
        hu_interpretable_2026,
        schlegel_towards_2025,
        lin_shap-based_2025,
        argov_spex_2025,
        de_oliveira_explainable_2025,
        hamdi_spatiotemporal_2022,
        ribeiro_anchors_2018,
        hwang_xclusters_2023,
        portugal_framework_2024,
        cakmak_spatio-temporal_2021,
        gabidolla_optimal_2022,
        salih_perspective_2025,
        birant_st-dbscan_2007,
        shi_spatiotemporal_2019,
        choi_mdst-dbscan_2021,
        kisilevich_spatio-temporal_nodate,
        dasgupta_explainable_nodate,
        lundberg_local_2020,
        paparrizos_time-series_2025,
        sabbatini_exact_nodate,
        ribeiro_why_2016-1,
        lundberg_unified_2017,
        elamouri_constrained_2023,
        ye_time_2009,
        bonifati_time2feat_2022}

\bibliography{example_paper}
\bibliographystyle{icml2025}

\end{document}